\chapter{Getting Started}
\label{ch:getting-started}

\section{Requirements}

\Morendo\ needs a Java Development Kit, version 21 or later, on the
\fn{PATH}. Nothing else has to be installed: the build downloads the few
libraries it uses (Log4j, and for the optional modules Jackson, Jakarta
Messaging, Jakarta Servlet and JLine) from Maven Central the first time it
runs. Linux, macOS and Windows are all supported; the examples in this
manual use a Unix shell.

\section{Building from source}

The build is driven by \fn{bld}, a small tool checked into the repository,
whose whole build definition is the Java file \file{builder/Tasks.java}.
From the root of a checkout:

\begin{lstlisting}[style=shell]
./bld build        # download libraries, generate the parser, compile
./bld test         # also compile and run the test suite
./bld dist         # target/morendo-3.0.0.zip, a self-contained distribution
./bld list-tasks   # every task, with a one-line description
\end{lstlisting}

On Windows use \fn{bld.cmd} instead of \fn{./bld}. Everything the build
produces goes under \file{target/}; the downloaded libraries go under
\file{libs/}. Both directories can be deleted at any time and are recreated
by the next build.

\section{Running the shell}

After \fn{./bld build}, the \fn{morendo} script at the root of the checkout
starts the engine with the compiled classes on the class path:

\begin{lstlisting}[style=shell]
./morendo -shell   # interactive shell
./morendo -gui     # Swing graphical interface
\end{lstlisting}

Both flags may be given together; Chapter~\ref{ch:shell-gui} describes the
graphical interface. The shell prints a banner and a prompt;
each expression you type is evaluated and its value printed:

\begin{lstlisting}[style=shell]
$ ./morendo -shell
Copyright Jamocha Project http://sourceforge.net/projects/jamocha
Morendo Version 3.0.0
Morendo> (+ 1 2)
3
Morendo> (batch samples/only/only_1.clp)
true
Morendo> (facts)
f-1 (_initialFact)
f-2 (goal (emergency TRUE) (people "nil") (location "nil") )
f-3 (hero (name "superman") (status "occupied") (long "nil") (lat "nil") )
f-4 (hero (name "spiderman") (status "unoccupied") (long "nil") (lat "nil") )
for a total of 4
Morendo> (fire)
spiderman save the day
1
Morendo> (exit)
\end{lstlisting}

\fn{batch} loads a file of CLIPS constructs and commands, here a small
sample shipped with \Morendo; \fn{facts} lists working memory (fact
\fn{f-1} is the initial fact the engine always keeps); \fn{fire} runs the
rules that match and returns how many fired; \fn{exit} closes the engine
and ends the shell. File names are relative to the directory you started
\Morendo\ in, so run it from the root of the checkout to use the samples
as written.

\begin{note}
The function that runs the agenda is \fn{fire}, not \fn{run} as in CLIPS.
An expression that calls a function \Morendo\ does not know simply
evaluates to \fn{false}; it is not reported as an error. If a command seems
to do nothing, check its name in Appendix~\ref{app:functions}.
\end{note}

The shell has line editing and history (it uses JLine), and an expression
may span several lines: input is collected until the parentheses balance.
Press Ctrl-D at the prompt, or type \fn{(exit)}, to leave.

\section{The distribution archive}

\fn{./bld dist} produces \file{target/morendo-\version.zip}. Unpacked, it
contains

\begin{center}
\begin{tabular}{ll}
\toprule
\file{libs/} & the six \Morendo\ jars and the third-party libraries \\
\file{morendo}, \file{morendo.cmd} & launchers for Unix and Windows \\
\file{samples/} & the example rule files used throughout this manual \\
\file{benchmark/} & the Manners benchmark rules and data sets \\
\file{README.md}, \file{LICENSE}, \file{licenses/} & documentation and licences \\
\bottomrule
\end{tabular}
\end{center}

The launcher puts every jar in \file{libs/} on the class path, so
\fn{./morendo -shell} works the same way from an unpacked distribution as
from a checkout. The jars are split by module so that a program embedding
the engine takes only what it needs:

\begin{center}
\begin{tabularx}{\textwidth}{lL}
\toprule
\file{morendo-core} & the engine, the parser and every built-in function; needs only Log4j \\
\file{morendo-shell} & the \fn{morendo} entry point and the interactive shell (JLine) \\
\file{morendo-gui} & the Swing interface and the network viewer; adds the \fn{view} function \\
\file{morendo-service} & the embedding service and servlet integration (Jackson, Jakarta Servlet) \\
\file{morendo-messaging} & the Jakarta Messaging client and agent functions \\
\file{morendo-examples} & the example beans used by the samples \\
\bottomrule
\end{tabularx}
\end{center}

Optional jars register their functions with the engine automatically when
they are on the class path. Chapter~\ref{ch:embedding} covers embedding.

\section{Where things are}

\begin{center}
\begin{tabularx}{\textwidth}{lL}
\toprule
\file{samples/} & small, self-contained rule files: \file{only/}, \file{exists/}, \file{time/}, \file{defquery/}, \file{graphquery/}, \file{molap/}, \file{ruleset/} \\
\file{samples/configuration/} & JSON and XML configurations for the service module \\
\file{benchmark/manners/} & the classic Manners benchmark with 8 to 256 guests \\
\file{doc/} & design papers on the MOLAP cube, modification logic and the temporal extension, plus the diagrams used in Chapter~\ref{ch:analysis} \\
\file{src/core/javacc/clips.jj} & the grammar; Appendix~\ref{app:grammar} summarizes it \\
\bottomrule
\end{tabularx}
\end{center}
